% Part: first-order-logic
% Chapter: proof-systems
% Section: introduction

\documentclass[../../../include/open-logic-section]{subfiles}

\begin{document}

\iftag{FOL}
      {\olfileid{fol}{prf}{int}}
      {\olfileid{pl}{prf}{int}}

\olsection{పరిచయం}

తర్కవ్యవస్థలకు సాధారణంగా అర్థవిచారమూ, ఒక \tetoken{వ్యుత్పత్తి}{!!a{derivation}}
వ్యవస్థా ఉంటాయి. అర్థవిచారం సత్యం, సంతృప్తిపరచదగినతనం,
చెల్లుబాటుతనం, అర్థపర అనుగమనం వంటి భావనలకు సంబంధించినది.
అనుగమనాన్నీ, చెల్లుబాటుతనాన్నీ స్థాపించడానికి పూర్తిగా
వాక్యనిర్మాణాత్మకమైన పద్ధతిని అందించడం \tetoken{వ్యుత్పత్తి}{!!{derivation}} వ్యవస్థల
ఉద్దేశం. అలాంటి వ్యవస్థలో ఒక \tetoken{వ్యుత్పత్తి}{!!{derivation}} అనేది సాధారణంగా
\tetoken{వాక్యాలు}{!!{sentence}s} లేదా \tetoken{సూత్రాలతో}{!!{formula}s} ఏర్పడిన క్రమం (లేదా మరేదైనా
పరిమిత అమరిక) వంటి పరిమిత వాక్యనిర్మాణ వస్తువు కనుక అవి పూర్తిగా
వాక్యనిర్మాణాత్మకమైనవి. మంచి \tetoken{వ్యుత్పత్తి}{!!{derivation}} వ్యవస్థలో, ఇచ్చిన ఏ
\tetoken{వాక్యాలు}{!!{sentence}s} లేదా \tetoken{సూత్రాలు}{!!{formula}s} క్రమం లేదా అమరిక అయినా “సరైనదో”
కాదో యాంత్రికంగా తనిఖీ చేయగలగాలి.

మొదటిస్థాయి విధేయ తర్కానికి అత్యంత సరళమైన (చారిత్రకంగా మొదటివైన)
\tetoken{వ్యుత్పత్తి}{!!{derivation}} వ్యవస్థలు
\emph{స్వీకృతాధారితమైనవి}. అలాంటి వ్యవస్థలోని \tetoken{సూత్రాలు}{!!{formula}s} క్రమం
ఒక \tetoken{వ్యుత్పత్తి}{!!a{derivation}} కావాలంటే, అందులోని ప్రతి ఒక్క \tetoken{సూత్రం}{!!{formula}}
స్థిరమైన “స్వీకృతాల” సమితిలో ఉండాలి, లేదా క్రమంలో తనకంటే ముందున్న
\tetoken{సూత్రాలు}{!!{formula}s} నుంచి స్థిరమైన “నిగమన నియమాల” సమితిలోని ఒక నియమం
ద్వారా రావాలి. అంతేకాక, ఒక \tetoken{సూత్రం}{!!a{formula}} స్వీకృతమా, ఇతర
\tetoken{సూత్రాలు}{!!{formula}s} నుంచి ఏదైనా నిగమన నియమం ద్వారా సరిగా వచ్చిందా అనేది
యాంత్రికంగా తనిఖీ చేయగలగాలి. స్వీకృతాధారిత \tetoken{వ్యుత్పత్తి}{!!{derivation}} వ్యవస్థలను
వర్ణించడం, అధిసిద్ధాంతపరంగా నిర్వహించడం సులభం; కానీ వాటిలోని
\tetoken{వ్యుత్పత్తులను}{!!{derivation}s} చదివి అర్థం చేసుకోవడం, నిర్మించడం కష్టం.

\tetoken{వ్యుత్పత్తులను}{!!{derivation}s} నిర్మించడమో, పూర్తయిన \tetoken{వ్యుత్పత్తులను}{!!{derivation}s} అర్థం చేసుకోవడమో
సులభం చేయాలనే లక్ష్యంతో ఇతర \tetoken{వ్యుత్పత్తి}{!!{derivation}} వ్యవస్థలు అభివృద్ధి
చేయబడ్డాయి. సహజ నిగమనం, సత్య వృక్షాలు అని కూడా పిలిచే టాబ్లో
నిరూపణలు, సీక్వెంట్ కలనం ఉదాహరణలు. కొన్ని \tetoken{వ్యుత్పత్తి}{!!{derivation}} వ్యవస్థలు
ముఖ్యంగా యాంత్రీకరణను దృష్టిలో పెట్టుకొని రూపొందించబడ్డాయి;
ఉదాహరణకు, రిజల్యూషన్ పద్ధతిని సాఫ్ట్‌వేర్‌లో అమలు చేయడం సులభం
(కానీ దాని \tetoken{వ్యుత్పత్తులను}{!!{derivation}s} అర్థం చేసుకోవడం దాదాపు అసాధ్యం). ఈ
ఇతర \tetoken{వ్యుత్పత్తి}{!!{derivation}} వ్యవస్థల్లో చాలా, \tetoken{వ్యుత్పత్తులను}{!!{derivation}s} క్రమాలుగా
కాకుండా \tetoken{సూత్రాలు}{!!{formula}s} వృక్షాలుగా సూచిస్తాయి. దీనివల్ల ఒక
\tetoken{వ్యుత్పత్తిలోని}{!!a{derivation}} ఏ భాగాలు మరే భాగాలపై ఆధారపడుతున్నాయో చూడడం
సులభమవుతుంది.

కాబట్టి మొదటిస్థాయి విధేయ తర్కం వంటి ఒక తర్కవ్యవస్థకు, వేర్వేరు
\tetoken{వ్యుత్పత్తి}{!!{derivation}} వ్యవస్థలు ఒక \tetoken{వాక్యం}{!!a{sentence}}
\emph{సిద్ధాంతం} కావడం అంటే ఏమిటో, కొన్ని ఇతర వాక్యాల నుంచి
ఒక \tetoken{వాక్యం}{!!a{sentence}} \tetoken{వ్యుత్పాదించదగిన}{!!{derivable}} కావడం అంటే ఏమిటో వేర్వేరుగా
వివరిస్తాయి. అది ఎలా చేసినా (స్వీకృతాధారిత \tetoken{వ్యుత్పత్తులు}{!!{derivation}s}, సహజ
నిగమనాలు, సీక్వెంట్ \tetoken{వ్యుత్పత్తులు}{!!{derivation}s}, సత్య వృక్షాలు, రిజల్యూషన్
నిరాకరణల ద్వారా), ఈ సంబంధాలు చెల్లుబాటుతనం, అర్థపర అనుగమనం అనే
అర్థపర భావనలతో సరిపోవాలని కోరుకుంటాం. “$!A$ ఒక సిద్ధాంతం” అని
సూచించడానికి $\Proves !A$ అనీ, “$\Gamma$ నుంచి $!A$
\tetoken{వ్యుత్పాదించదగిన}{!!{derivable}}” అని సూచించడానికి “$\Gamma \Proves !A$” అనీ రాద్దాం.
$\Proves$ను ఎలా నిర్వచించినా అది $\Entails$తో సరిపోవాలని, అంటే
కిందివి కావాలని కోరుకుంటాం:
\begin{enumerate}
\item $\Proves !A$ కావడం $\Entails !A$ అయినప్పుడూ, అప్పుడు మాత్రమే
\item $\Gamma \Proves !A$ కావడం $\Gamma \Entails !A$ అయినప్పుడూ,
  అప్పుడు మాత్రమే
\end{enumerate}
పై సమానార్థక ప్రకటనల్లో నిగమ్యత నుంచి అర్థపర అనుగమనానికి వెళ్లే
“అప్పుడు మాత్రమే” దిశను \emph{నిర్దుష్టత} అంటారు.
\tetoken{ఒక వ్యుత్పత్తి}{!!^a{derivation}} వ్యవస్థలో \tetoken{వ్యుత్పాద్యత}{!!{derivability}} వల్ల అనుగమనం (లేదా
చెల్లుబాటుతనం) తప్పనిసరిగా వస్తే ఆ వ్యవస్థ నిర్దుష్టమైనది. ఉపయోగకరమైన
ప్రతి \tetoken{వ్యుత్పత్తి}{!!{derivation}} వ్యవస్థా నిర్దుష్టంగా ఉండాలి; నిర్దుష్టం కాని
\tetoken{వ్యుత్పత్తి}{!!{derivation}} వ్యవస్థలు అసలు ఉపయోగపడవు. ఎందుకంటే చెల్లుబాటుతనం
లేదా అనుగమనానికి వాక్యనిర్మాణాత్మక హామీ ఇవ్వడమే ఒక
\tetoken{వ్యుత్పత్తి}{!!a{derivation}} యొక్క పూర్తి ఉద్దేశం. మనం సమర్పించే \tetoken{వ్యుత్పత్తి}{!!{derivation}}
వ్యవస్థల నిర్దుష్టతను నిరూపిస్తాం.

దీనికి విరుద్ధమైన “అప్పుడు” దిశ కూడా ముఖ్యం; దానిని
\emph{సంపూర్ణత} అంటారు. $!A$ చెల్లుబాటయ్యే ప్రతిసారీ $!A$ ఒక
సిద్ధాంతమని చూపగలిగేంత, అలాగే $\Gamma \Entails !A$ అయినప్పుడల్లా
$\Gamma \Proves !A$ను చూపగలిగేంత బలమున్న \tetoken{వ్యుత్పత్తి}{!!{derivation}} వ్యవస్థ
సంపూర్ణమైనది. సంపూర్ణతను స్థాపించడం మరింత కష్టం; కొన్ని
తర్కవ్యవస్థలకు సంపూర్ణమైన \tetoken{వ్యుత్పత్తి}{!!{derivation}} వ్యవస్థలు లేవు. మొదటిస్థాయి
తర్కానికి ఉన్నాయి. కుర్ట్ గోడెల్ తన 1929 పరిశోధన గ్రంథంలో మొదటిస్థాయి
తర్కానికి ఒక \tetoken{వ్యుత్పత్తి}{!!a{derivation}} వ్యవస్థ సంపూర్ణమని మొదట నిరూపించాడు.

\tetoken{వ్యుత్పత్తి}{!!{derivation}} వ్యవస్థలకు సంబంధించిన మరొక భావన
\emph{అవైరుధ్యం}. ఒక \tetoken{వాక్యాలు}{!!{sentence}s} సమితి నుంచి ఏదైనా సరే
\tetoken{నిగమనం}{!!{derive}} చేయగలిగితే దానిని వైరుధ్యభరితమనీ, లేకపోతే అవైరుధ్యమైనదనీ
అంటారు. వైరుధ్యం అనేది అసంతృప్తిపరచలేనితనానికి వాక్యనిర్మాణాత్మక
ప్రతిరూపం: సంతృప్తిపరచలేని సమితుల మాదిరిగానే, వైరుధ్యభరితమైన
\tetoken{వాక్యాలు}{!!{sentence}s} సమితులు మంచి సిద్ధాంతాలు కావు; వాటిలో మౌలిక లోపం
ఉంటుంది. అవైరుధ్యమైన \tetoken{వాక్యాలు}{!!{sentence}s} సమితులు సత్యమైనవో ఉపయోగకరమైనవో
కాకపోవచ్చు, కానీ కనీసం తార్కిక ఉపయోగానికి అవసరమైన ఈ కనిష్ఠ ప్రమాణాన్ని
దాటుతాయి. వేర్వేరు \tetoken{వ్యుత్పత్తి}{!!{derivation}} వ్యవస్థల్లో \tetoken{వాక్యాలు}{!!{sentence}s} సమితుల
అవైరుధ్యానికి నిర్దిష్ట నిర్వచనం మారవచ్చు; కానీ $\Proves$ మాదిరిగానే,
అవైరుధ్యం దాని అర్థపర ప్రతిరూపమైన సంతృప్తిపరచదగినతనంతో సరిపోవాలని
కోరుకుంటాం. $\Gamma$ అవైరుధ్యంగా ఉండడం అది సంతృప్తిపరచదగినప్పుడు,
అప్పుడు మాత్రమే కావాలి. ఇక్కడ అవైరుధ్యం నుంచి
సంతృప్తిపరచదగినతనానికి వెళ్లే “అప్పుడు మాత్రమే” దిశ సంపూర్ణతకు,
సంతృప్తిపరచదగినతనం నుంచి అవైరుధ్యానికి వెళ్లే తిరుగు దిశ
నిర్దుష్టతకు సమానం. వాస్తవానికి, సాంప్రదాయ మొదటిస్థాయి విధేయ తర్కంలో
నిర్దుష్టత, సంపూర్ణతల రెండు రూపాలు సమానార్థకాలు.

\end{document}
